% 自旋求和、偏振密度矩阵与外部实光子偏振代数

树级费米子过程反复需要把未分辨外部自旋转成 Dirac 矩阵迹。物理四动量统一取
\begin{equation*}
 p^\mu=(E_{\mathbf p}/c,\mathbf p),
 \qquad p^2=m^2c^2,
\end{equation*}
并使用精细结构常数
\begin{equation*}
 \alpha\equiv\frac{e^2}{4\pi\varepsilon_0\hbar c},
 \qquad e>0,
 \qquad q_{\mathrm e}=-e.
\end{equation*}
因此所有自旋子完备关系中的质量项都必须写成动量量纲的 \(mc\)。

正、负能自由旋量的协变完备关系为
\begin{equation}
 \boxed{
 \sum_su_s(p)\bar u_s(p)=\slashed p+mc,
 \qquad
 \sum_sv_s(p)\bar v_s(p)=\slashed p-mc.}
 \label{eq:spin-sum-si-r9}
\end{equation}
结合 Clifford 迹恒等式，矢量流中最常复用的秩二张量直接化为
\begin{equation}
 \boxed{
 L^{\mu\nu}(a,b;m)
 \equiv
 \Tr[(\slashed a+mc)\gamma^\mu(\slashed b+mc)\gamma^\nu]
 =4\!
 \left[a^\mu b^\nu+a^\nu b^\mu
 -(a\!\cdot b-m^2c^2)\eta^{\mu\nu}\right].}
 \label{eq:qed-basic-tensor-si-r9}
\end{equation}

\begin{proposition}[Dirac 自旋求和的迹恒等式]
设外态在壳，gamma 矩阵满足正文 Minkowski Clifford 关系，采用正文自旋态归一化。则有以下结论（原式 \eqnrefs{eq:qed-basic-tensor-si-r9}）：
\begin{equation*}
 \boxed{
 L^{\mu\nu}(a,b;m)
 \equiv
 \Tr[(\slashed a+mc)\gamma^\mu(\slashed b+mc)\gamma^\nu]
 =4\!
 \left[a^\mu b^\nu+a^\nu b^\mu
 -(a\!\cdot b-m^2c^2)\eta^{\mu\nu}\right].}
\end{equation*}
\end{proposition}
\begin{proof}
未偏振截面需要对末态自旋求和、对初态自旋平均。正文\eqnrefs{eq:spin-sum-si-r9} 把这些离散自旋指标变成 $\slashed p+mc$，随后利用迹循环性可把两条外部自旋链合并成单个 Dirac 迹；对矢量流执行这一步便得到正文\eqnrefs{eq:qed-basic-tensor-si-r9}。



取
\begin{equation*}
 B_{rs}=\bar u_r(p')\mathcal O u_s(p),
\end{equation*}
其中 $\mathcal O$ 是任意 $4\times4$ 旋量空间矩阵；这里的花体 $\mathcal O$ 只表示插在两个外线自旋子之间的算符，不与 1PI 顶角 $\Gamma^\mu$ 或衰变率记号混用。先求复共轭：
\begin{align*}
 B_{rs}^*
 &=[u_r^\dagger(p')\gamma^0\mathcal O u_s(p)]^*,\\
 &=u_s^\dagger(p)\mathcal O^\dagger\gamma^0u_r(p'),\\
 &=\bar u_s(p)\gamma^0\mathcal O^\dagger\gamma^0u_r(p').
\end{align*}
定义
\begin{equation*}
 \overline{\mathcal O}\equiv\gamma^0\mathcal O^\dagger\gamma^0.
\end{equation*}
于是
\begin{align*}
 \sum_{r,s}|B_{rs}|^2
 &=\sum_{r,s}
 \bar u_r(p')\mathcal O u_s(p)
 \bar u_s(p)\overline{\mathcal O} u_r(p'),\\
 &=\sum_r\bar u_r(p')\mathcal O
 \left[\sum_su_s(p)\bar u_s(p)\right]
 \overline{\mathcal O} u_r(p'),\\
 &=\sum_r\bar u_r(p')\mathcal O(\slashed p+mc)
 \overline{\mathcal O} u_r(p').
\end{align*}
把最后一个标量写成迹：
\begin{equation*}
 \bar u_r X u_r=\Tr[Xu_r\bar u_r],
\end{equation*}
再对 $r$ 求和，
\begin{align*}
 \sum_{r,s}|B_{rs}|^2
 &=\Tr\!\left[
 \mathcal O(\slashed p+mc)\overline{\mathcal O}
 \sum_ru_r(p')\bar u_r(p')
 \right],\\
 &={
 \Tr\!\left[(\slashed p'+mc)\mathcal O
 (\slashed p+mc)\overline{\mathcal O}\right].}
 \end{align*}
未偏振入射态由自旋密度矩阵 $\mathbb I_2/2$ 表示，末态自旋求和由完备关系实现；两者把有限个外线自旋指标精确收缩成矩阵迹。

对矢量流取 $\mathcal O=\gamma^\mu$，并使用
\begin{equation*}
 \Tr(\gamma^\alpha\gamma^\beta)=4\eta^{\alpha\beta},
 \qquad
 \Tr(\gamma^\alpha\gamma^\mu\gamma^\beta\gamma^\nu)
 =4(\eta^{\alpha\mu}\eta^{\beta\nu}
 -\eta^{\alpha\beta}\eta^{\mu\nu}
 +\eta^{\alpha\nu}\eta^{\mu\beta}),
\end{equation*}
奇数个 $\gamma$ 矩阵的迹为零，于是
\begin{align*}
 \Tr[(\slashed a+mc)\gamma^\mu(\slashed b+mc)\gamma^\nu]
 &=\Tr(\slashed a\gamma^\mu\slashed b\gamma^\nu)
 +m^2c^2\Tr(\gamma^\mu\gamma^\nu),\\
 &=4[a^\mu b^\nu+a^\nu b^\mu-(a\!\cdot b-m^2c^2)\eta^{\mu\nu}],
\end{align*}
即\eqnrefs{eq:qed-basic-tensor-si-r9}。
\end{proof}

未偏振入射态由自旋密度矩阵 $\rho_s=\mathbb I_2/2$ 定义，末态自旋不分辨则由完备求和实现。将这两个统计操作与\eqnrefs{eq:spin-sum-si-r9} 结合后，矩阵元平方化为相应的 Dirac 迹张量。

若实验准备了确定自旋方向，令 \(\boldsymbol\zeta\) 为粒子静止系中的单位自旋方向，\(|\boldsymbol\zeta|=1\)。把 \((0,\boldsymbol\zeta)\) 与粒子四动量用同一个 Lorentz 推动到实验参考系，得到自旋四矢量
\begin{equation}
 \boxed{
 s^0=\gamma\frac{\mathbf v\!\cdot\boldsymbol\zeta}{c},
 \qquad
 \mathbf s
 =\boldsymbol\zeta
 +\frac{\gamma^2}{\gamma+1}
 \frac{\mathbf v(\mathbf v\!\cdot\boldsymbol\zeta)}{c^2},}
 \label{eq:spin-fourvector-r9}
\end{equation}
满足 \(s^2=-1\) 与 \(s\!\cdot p=0\)。对应的固定自旋密度矩阵为
\begin{equation}
 \boxed{
 u(p,s)\bar u(p,s)
 =\frac12(\slashed p+mc)(1+\gamma^5\slashed s).}
 \label{eq:polarized-spinor-projector-r9}
\end{equation}

\begin{proposition}[极化 Dirac 态的投影矩阵]
设正能四动量与自旋四矢量满足 $p^2=m^2c^2$、$s^2=-1$、$p\cdot s=0$，并固定旋量归一化。则有以下结论（原式 \eqnrefs{eq:polarized-spinor-projector-r9}）：
\begin{equation*}
 \boxed{
 u(p,s)\bar u(p,s)
 =\frac12(\slashed p+mc)(1+\gamma^5\slashed s).}
\end{equation*}
\end{proposition}
\begin{proof}
Landau--Lifshitz 型的偏振散射计算经常不先做自旋求和，而是保留一个指定的极化 four-矢量。先在粒子静止系取单位三矢量 $\boldsymbol\zeta$ 指定自旋方向，定义

\begin{equation*}
 s_{\rm rest}^\mu=(0,\boldsymbol\zeta),
 \qquad
 s_{\rm rest}^2=-1,
 \qquad
 s_{\rm rest}\!\cdot p_{\rm rest}=0.
\end{equation*}
把 $p_{\rm rest}^\mu=(mc,\mathbf0)$ 与 $s_{\rm rest}^\mu$ 用同一个 Lorentz 推动变换到一般参考系。若 $\mathbf v$ 是粒子速度、$\gamma=E/(mc^2)$，推动给
\begin{equation*}
 {
 s^0=\gamma\frac{\mathbf v\!\cdot\boldsymbol\zeta}{c},
 \qquad
 \mathbf s
 =\boldsymbol\zeta
 +\frac{\gamma^2}{\gamma+1}
 \frac{\mathbf v(\mathbf v\!\cdot\boldsymbol\zeta)}{c^2}.}
 \end{equation*}
直接代入可检查
\begin{equation*}
 s^2=-1,
 \qquad
 s\!\cdot p=0.
\end{equation*}
在静止系，正能旋量的自旋投影算符是
\begin{equation*}
 \frac12(1+\boldsymbol\Sigma\!\cdot\boldsymbol\zeta).
\end{equation*}
利用
\begin{equation*}
 \gamma^5\slashed s_{\rm rest}
 =\boldsymbol\Sigma\!\cdot\boldsymbol\zeta
\end{equation*}
和正能投影算符 $\slashed p_{\rm rest}+mc=mc(\gamma^0+1)$，单一自旋态满足
\begin{equation*}
 u(p_{\rm rest},s)\bar u(p_{\rm rest},s)
 =\frac12(\slashed p_{\rm rest}+mc)
 (1+\gamma^5\slashed s_{\rm rest}).
\end{equation*}
两边在 Lorentz 变换下都按 $S(\Lambda)(\cdots)S^{-1}(\Lambda)$ 变换，所以一般参考系中
\begin{equation*}
 {
 u(p,s)\bar u(p,s)
 =\frac12(\slashed p+mc)(1+\gamma^5\slashed s).}
 \end{equation*}
把相反偏振 $s^\mu\to-s^\mu$ 的两式相加，$\gamma^5\slashed s$ 项消失，所得和等于正文的未偏振完备关系。
\end{proof}

对相反的两个自旋方向求和即恢复\eqnrefs{eq:spin-sum-si-r9}。

外部实光子的两个物理自由度已经由\eqnrefs{eq:physical-photon-polarization-projector-dp42} 的横向完备张量定义。若振幅写成 \(\mathcal M=\varepsilon_\mu\mathcal M^\mu\)，并满足 Ward 条件 \(k_\mu\mathcal M^\mu=0\)，则辅助参考矢量相关项全部消失，\eqnrefs{eq:physical-photon-polarization-conserved-current-dp42} 允许在\emph{与守恒振幅完成收缩以后}使用
\begin{equation*}
 \sum_{\lambda=1,2}
 \varepsilon_\lambda^\mu\varepsilon_\lambda^{\nu *}
 \rightsquigarrow-\eta^{\mu\nu}.
\end{equation*}
这个箭头不是完整四维张量恒等式，只表示物理横向偏振求和在守恒矩阵元中的等效替换。

作用量与顶角先确定量子振幅；未分辨的自旋或光子偏振只在振幅构造完成后通过密度矩阵和物理偏振投影进入 $\overline{|\mathcal M|^2}$，再与相空间和入射通量组成实验截面。因而自旋平均、固定自旋和偏振求和属于外态制备与读取规则，不改变相互作用 Hamiltonian。

